% GENERATED FILE. Edit canonical lessons, then run npm run generate:books.
% canonical-source-hash: fnv1a64:08c1fa4566a9635d
% canonical-lessons: ML-C404-caritrakaran, ML-C404-puratana, ML-C404-taliyola, ML-C404-lipi, ML-C404-nasippikkuka, ML-R404-review-397-404, ML-R404-review-final-404

\chapter{Old Writings and Ruins}
\label{ch:ml-old-writings-and-ruins}
\hlchaptermodality{404}

\begin{chapteropening}
\textbf{By the end of this chapter:} \emph{I can say a historian, ancient, a palm-leaf manuscript, a script and to destroy.}

Complete the last lesson of chapter 404: I can say a historian, ancient, a palm-leaf manuscript, a script and to destroy.
\end{chapteropening}

\section[\texorpdfstring{\ml{ചരിത്രകാരൻ} (caritrakāran)}{ചരിത്രകാരൻ (caritrakāran)}]{\ml{ചരിത്രകാരൻ} (caritrakāran) — a historian}

\label{lesson:ML-C404-caritrakaran}



\begin{quote}
Before the new one: say the Malayalam for a curse, then the Malayalam for to curse.
\end{quote}

\subsection*{You'll want to know: \ml{ചരിത്രകാരൻ}}
\textbf{\ml{ചരിത്രകാരൻ}} — \emph{caritrakāran} — \textquotedblleft{}a historian\textquotedblright{}.

\begin{quote}
\ml{ചരിത്രകാരൻ} \ml{ആ} \ml{കോട്ടയെക്കുറിച്ച്} \ml{ഒരു} \ml{പുസ്തകം} \ml{എഴുതി}.
\end{quote}

\begin{quote}
\emph{caritrakāran ā kōṭṭayekkuṟiccŭ oru pustakaṁ eḻuti.}
\end{quote}

\begin{quote}
The historian wrote a book about that fort.
\end{quote}

\ml{ചരിത്രം}, history, + -\ml{കാരൻ}, one who works at it.

\begin{cousinweb}
\textbf{\ml{ചരിത്രകാരൻ}} (\emph{caritrakāran}) is built on \ml{ചരിത്രം} (\emph{caritraṁ}), \textquotedblleft{}history\textquotedblright{}, from Sanskrit \emph{caritra}, \textquotedblleft{}going; behaviour\textquotedblright{} (CDIAL 4691), with the ending -\ml{കാരൻ}.
\end{cousinweb}

\begin{grammarlens}[title={What you've built}]
One more word for telling what happened.
\end{grammarlens}

\subsection*{Your turn}
\begin{itemize}
\raggedright
  \item \emph{Say it:} \emph{caritrakāran}
  \item \emph{Say it:} \emph{caritrakāran}, once more
  \item \emph{Say it:} \emph{caritrakāran ā kōṭṭayekkuṟiccŭ oru pustakaṁ eḻuti.} — The historian wrote a book about that fort.
  \item \emph{Recall:} say the Malayalam for a curse, then the Malayalam for to curse, then say \emph{caritrakāran} again
\end{itemize}

\begin{tcolorbox}[breakable,colback=teal!4,colframe=teal!35!black,title={Before you move on}]
What does \ml{ചരിത്രകാരൻ} mean? (\textquotedblleft{}A historian\textquotedblright{}.) Say it once more.
\end{tcolorbox}
\section[\texorpdfstring{\ml{പുരാതന} (purātana)}{പുരാതന (purātana)}]{\ml{പുരാതന} (purātana) — ancient}

\label{lesson:ML-C404-puratana}



\begin{quote}
Before the new one: say the Malayalam for to curse, then the Malayalam for a historian.
\end{quote}

\subsection*{You'll want to know: \ml{പുരാതന}}
\textbf{\ml{പുരാതന}} — \emph{purātana} — \textquotedblleft{}ancient\textquotedblright{}.

\begin{quote}
\ml{മ്യൂസിയത്തിൽ} \ml{ഒരുപാട്} \ml{പുരാതന} \ml{നാണയങ്ങൾ} \ml{ഉണ്ട്}.
\end{quote}

\begin{quote}
\emph{myūsiyattil orupāṭŭ purātana nāṇayaṅṅaḷ uṇṭŭ.}
\end{quote}

\begin{quote}
There are many ancient coins in the museum.
\end{quote}

It stands before the noun, unchanged: \ml{പുരാതന} \ml{ക്ഷേത്രം}, an ancient temple.

\begin{cousinweb}
\textbf{\ml{പുരാതന}} (\emph{purātana}) is borrowed from Sanskrit \emph{purātana}, \textquotedblleft{}ancient, belonging to the past\textquotedblright{} (CDIAL 8285).
\end{cousinweb}

\begin{grammarlens}[title={What you've built}]
One more word for telling what happened.
\end{grammarlens}

\subsection*{Your turn}
\begin{itemize}
\raggedright
  \item \emph{Say it:} \emph{purātana}
  \item \emph{Say it:} \emph{purātana}, once more
  \item \emph{Say it:} \emph{myūsiyattil orupāṭŭ purātana nāṇayaṅṅaḷ uṇṭŭ.} — There are many ancient coins in the museum.
  \item \emph{Recall:} say the Malayalam for to curse, then the Malayalam for a historian, then say \emph{purātana} again
\end{itemize}

\begin{tcolorbox}[breakable,colback=teal!4,colframe=teal!35!black,title={Before you move on}]
What does \ml{പുരാതന} mean? (\textquotedblleft{}Ancient\textquotedblright{}.) Say it once more.
\end{tcolorbox}
\section[\texorpdfstring{\ml{താളിയോല} (tāḷiyōla)}{താളിയോല (tāḷiyōla)}]{\ml{താളിയോല} (tāḷiyōla) — a palm-leaf manuscript}

\label{lesson:ML-C404-taliyola}



\begin{quote}
Before the new one: say the Malayalam for a historian, then the Malayalam for ancient.
\end{quote}

\subsection*{You'll want to know: \ml{താളിയോല}}
\textbf{\ml{താളിയോല}} — \emph{tāḷiyōla} — \textquotedblleft{}a palm-leaf manuscript\textquotedblright{}.

\begin{quote}
\ml{ഗ്രന്ഥശാലയിൽ} \ml{പഴയ} \ml{താളിയോലകൾ} \ml{സൂക്ഷിച്ചിരിക്കുന്നു}.
\end{quote}

\begin{quote}
\emph{granthaśālayil paḻaya tāḷiyōlakaḷ sūkṣiccirikkunnu.}
\end{quote}

\begin{quote}
Old palm-leaf manuscripts are kept in the library.
\end{quote}

Plural: \ml{താളിയോലകൾ}, \emph{tāḷiyōlakaḷ}.

Before paper, books were written on such leaves.

\begin{cousinweb}
\textbf{\ml{താളിയോല}} (\emph{tāḷiyōla}) is a native compound of \ml{താളി}, usually traced to the name of a palm, and \ml{ഓല} (\emph{ōla}), \textquotedblleft{}a palm leaf, a writing leaf\textquotedblright{} (DEDR 1070).
\end{cousinweb}

\begin{grammarlens}[title={What you've built}]
One more word for telling what happened.
\end{grammarlens}

\subsection*{Your turn}
\begin{itemize}
\raggedright
  \item \emph{Say it:} \emph{tāḷiyōla}
  \item \emph{Say it:} \emph{tāḷiyōla}, once more
  \item \emph{Say it:} \emph{granthaśālayil paḻaya tāḷiyōlakaḷ sūkṣiccirikkunnu.} — Old palm-leaf manuscripts are kept in the library.
  \item \emph{Recall:} say the Malayalam for a historian, then the Malayalam for ancient, then say \emph{tāḷiyōla} again
\end{itemize}

\begin{tcolorbox}[breakable,colback=teal!4,colframe=teal!35!black,title={Before you move on}]
What does \ml{താളിയോല} mean? (\textquotedblleft{}A palm-leaf manuscript\textquotedblright{}.) Say it once more.
\end{tcolorbox}
\section[\texorpdfstring{\ml{ലിപി} (lipi)}{ലിപി (lipi)}]{\ml{ലിപി} (lipi) — a script; writing}

\label{lesson:ML-C404-lipi}



\begin{quote}
Before the new one: say the Malayalam for ancient, then the Malayalam for a palm-leaf manuscript.
\end{quote}

\subsection*{You'll want to know: \ml{ലിപി}}
\textbf{\ml{ലിപി}} — \emph{lipi} — \textquotedblleft{}a script; writing\textquotedblright{}.

\begin{quote}
\ml{മലയാള} \ml{ലിപിയിൽ} \ml{ഒരുപാട്} \ml{അക്ഷരങ്ങൾ} \ml{ഉണ്ട്}.
\end{quote}

\begin{quote}
\emph{malayāḷa lipiyil orupāṭŭ akṣaraṅṅaḷ uṇṭŭ.}
\end{quote}

\begin{quote}
The Malayalam script has many letters.
\end{quote}

Also any system of writing.

\begin{cousinweb}
\textbf{\ml{ലിപി}} (\emph{lipi}) is borrowed from Sanskrit \emph{lipi}, \textquotedblleft{}writing\textquotedblright{} (CDIAL 11058); CDIAL traces the Sanskrit word itself to Old Persian \emph{dipi}, \textquotedblleft{}an inscription\textquotedblright{}.
\end{cousinweb}

\begin{grammarlens}[title={What you've built}]
One more word for telling what happened.
\end{grammarlens}

\subsection*{Your turn}
\begin{itemize}
\raggedright
  \item \emph{Say it:} \emph{lipi}
  \item \emph{Say it:} \emph{lipi}, once more
  \item \emph{Say it:} \emph{malayāḷa lipiyil orupāṭŭ akṣaraṅṅaḷ uṇṭŭ.} — The Malayalam script has many letters.
  \item \emph{Recall:} say the Malayalam for ancient, then the Malayalam for a palm-leaf manuscript, then say \emph{lipi} again
\end{itemize}

\begin{tcolorbox}[breakable,colback=teal!4,colframe=teal!35!black,title={Before you move on}]
What does \ml{ലിപി} mean? (\textquotedblleft{}A script; writing\textquotedblright{}.) Say it once more.
\end{tcolorbox}
\section[\texorpdfstring{\ml{നശിപ്പിക്കുക} (naśippikkuka)}{നശിപ്പിക്കുക (naśippikkuka)}]{\ml{നശിപ്പിക്കുക} (naśippikkuka) — to destroy}

\label{lesson:ML-C404-nasippikkuka}



\begin{quote}
Before the new one: say the Malayalam for a palm-leaf manuscript, then the Malayalam for a script.
\end{quote}

\subsection*{You'll want to know: \ml{നശിപ്പിക്കുക}}
\textbf{\ml{നശിപ്പിക്കുക}} — \emph{naśippikkuka} — \textquotedblleft{}to destroy\textquotedblright{}.

\begin{quote}
\ml{യുദ്ധം} \ml{ആ} \ml{പുരാതന} \ml{നഗരം} \ml{നശിപ്പിച്ചു}.
\end{quote}

\begin{quote}
\emph{yuddhaṁ ā purātana nagaraṁ naśippiccu.}
\end{quote}

\begin{quote}
The war destroyed that ancient city.
\end{quote}

Past: \ml{നശിപ്പിച്ചു}, \emph{naśippiccu}.

Said of war, of fire and of time: what is destroyed is the plain object.

\begin{cousinweb}
\textbf{\ml{നശിപ്പിക്കുക}} (\emph{naśippikkuka}) is built from Sanskrit \emph{nāśana}, \textquotedblleft{}destroying\textquotedblright{} (CDIAL 7086), with the Malayalam causative ending -\ml{ിപ്പിക്കുക}.
\end{cousinweb}

\begin{grammarlens}[title={What you've built}]
One more word for telling what happened.
\end{grammarlens}

\subsection*{Your turn}
\begin{itemize}
\raggedright
  \item \emph{Say it:} \emph{naśippikkuka}
  \item \emph{Say it:} \emph{naśippikkuka}, once more
  \item \emph{Say it:} \emph{yuddhaṁ ā purātana nagaraṁ naśippiccu.} — The war destroyed that ancient city.
  \item \emph{Recall:} say the Malayalam for a palm-leaf manuscript, then the Malayalam for a script, then say \emph{naśippikkuka} again
\end{itemize}

\begin{tcolorbox}[breakable,colback=teal!4,colframe=teal!35!black,title={Before you move on}]
What does \ml{നശിപ്പിക്കുക} mean? (\textquotedblleft{}To destroy\textquotedblright{}.) Say it once more.
\end{tcolorbox}
\section[(dialogue)]{Words for telling what happened, once more}

\label{lesson:ML-R404-review-397-404}



\begin{quote}
Say the Malayalam for a king and a queen.
\end{quote}

\subsection*{Your turn}
Say these aloud:

\begin{itemize}
\raggedright
  \item a king, a queen, a prince, a princess, a palace
  \item a throne, a crown, a coronation, a sovereign, to reign
  \item an empire, a fort, an army, a soldier, to conquer
  \item an enemy, a warrior, a spy, to beat in battle, to surrender
  \item a slave, a revolution, an invasion, to fight, to resist
  \item a myth, a fable, a folk tale, a hero, a moral
  \item a ghost, a magician, a treasure, a curse, to curse
  \item a historian, ancient, a palm-leaf manuscript, a script, to destroy
\end{itemize}

\begin{tcolorbox}[breakable,colback=teal!4,colframe=teal!35!black,title={Before you move on}]
A king? (\textbf{\ml{രാജാവ്}}, \emph{rājāvŭ}.) A queen? (\textbf{\ml{രാജ്ഞി}}, \emph{rājñi}.) A prince? (\textbf{\ml{രാജകുമാരൻ}}, \emph{rājakumāran}.) A princess? (\textbf{\ml{രാജകുമാരി}}, \emph{rājakumāri}.) A palace? (\textbf{\ml{കൊട്ടാരം}}, \emph{koṭṭāraṁ}.) A throne? (\textbf{\ml{സിംഹാസനം}}, \emph{siṁhāsanaṁ}.) A crown? (\textbf{\ml{കിരീടം}}, \emph{kirīṭaṁ}.) A coronation? (\textbf{\ml{പട്ടാഭിഷേകം}}, \emph{paṭṭābhiṣēkaṁ}.) A sovereign? (\textbf{\ml{ഭരണാധികാരി}}, \emph{bharaṇādhikāri}.) To reign? (\textbf{\ml{വാഴുക}}, \emph{vāḻuka}.) An empire? (\textbf{\ml{സാമ്രാജ്യം}}, \emph{sāmrājyaṁ}.) A fort? (\textbf{\ml{കോട്ട}}, \emph{kōṭṭa}.) An army? (\textbf{\ml{സൈന്യം}}, \emph{sainyaṁ}.) A soldier? (\textbf{\ml{പട്ടാളക്കാരൻ}}, \emph{paṭṭāḷakkāran}.) To conquer? (\textbf{\ml{കീഴടക്കുക}}, \emph{kīḻaṭakkuka}.) An enemy? (\textbf{\ml{ശത്രു}}, \emph{śatru}.) A warrior? (\textbf{\ml{യോദ്ധാവ്}}, \emph{yōddhāvŭ}.) A spy? (\textbf{\ml{ചാരൻ}}, \emph{cāran}.) To beat in battle? (\textbf{\ml{തോൽപ്പിക്കുക}}, \emph{tōlppikkuka}.) To surrender? (\textbf{\ml{കീഴടങ്ങുക}}, \emph{kīḻaṭaṅṅuka}.) A slave? (\textbf{\ml{അടിമ}}, \emph{aṭima}.) A revolution? (\textbf{\ml{വിപ്ലവം}}, \emph{viplavaṁ}.) An invasion? (\textbf{\ml{അധിനിവേശം}}, \emph{adhinivēśaṁ}.) To fight? (\textbf{\ml{പോരാടുക}}, \emph{pōrāṭuka}.) To resist? (\textbf{\ml{ചെറുത്തുനിൽക്കുക}}, \emph{ceṟuttunilkkuka}.) A myth? (\textbf{\ml{പുരാണം}}, \emph{purāṇaṁ}.) A fable? (\textbf{\ml{കെട്ടുകഥ}}, \emph{keṭṭukatha}.) A folk tale? (\textbf{\ml{നാടോടിക്കഥ}}, \emph{nāṭōṭikkatha}.) A hero? (\textbf{\ml{നായകൻ}}, \emph{nāyakan}.) A moral? (\textbf{\ml{ഗുണപാഠം}}, \emph{guṇapāṭhaṁ}.) A ghost? (\textbf{\ml{പ്രേതം}}, \emph{prētaṁ}.) A magician? (\textbf{\ml{മാന്ത്രികൻ}}, \emph{māntrikan}.) A treasure? (\textbf{\ml{നിധി}}, \emph{nidhi}.) A curse? (\textbf{\ml{ശാപം}}, \emph{śāpaṁ}.) To curse? (\textbf{\ml{ശപിക്കുക}}, \emph{śapikkuka}.) A historian? (\textbf{\ml{ചരിത്രകാരൻ}}, \emph{caritrakāran}.) Ancient? (\textbf{\ml{പുരാതന}}, \emph{purātana}.) A palm-leaf manuscript? (\textbf{\ml{താളിയോല}}, \emph{tāḷiyōla}.) A script? (\textbf{\ml{ലിപി}}, \emph{lipi}.) To destroy? (\textbf{\ml{നശിപ്പിക്കുക}}, \emph{naśippikkuka}.)
\end{tcolorbox}
\section[(dialogue)]{Old Writings and Ruins — once more}

\label{lesson:ML-R404-review-final-404}



\begin{quote}
Say the Malayalam for a historian and ancient.
\end{quote}

\subsection*{Your turn}
\begin{itemize}
\raggedright
  \item \emph{Say it:} a historian, ancient, a palm-leaf manuscript, a script, to destroy
\end{itemize}

\begin{tcolorbox}[breakable,colback=teal!4,colframe=teal!35!black,title={Before you move on}]
A historian? (\textbf{\ml{ചരിത്രകാരൻ}}, \emph{caritrakāran}.) Ancient? (\textbf{\ml{പുരാതന}}, \emph{purātana}.) A palm-leaf manuscript? (\textbf{\ml{താളിയോല}}, \emph{tāḷiyōla}.) A script? (\textbf{\ml{ലിപി}}, \emph{lipi}.) To destroy? (\textbf{\ml{നശിപ്പിക്കുക}}, \emph{naśippikkuka}.)
\end{tcolorbox}
